\documentclass{article}
\usepackage{propositions}

%% align=runin.  Like flush, in that the label sits at the text margin and no
%% label area is reserved, but the label is simply the first word of the body
%% and is followed by an ordinary interword space rather than \labelsep.  It
%% takes the same route as nextline and flush-nextline -- the label never
%% enters the label box -- which is why per-item dimensions have no bearing on
%% it.

\begin{document}

\section{\texttt{runin} beside \texttt{flush} and the default}

The three lists differ only in \texttt{align}.  Under \verb|runin| the gap
after the label must be one interword space, and the first line must justify
like any other; under \verb|flush| it is a \verb|\labelsep| and the first line
is set within the label machinery; under the default the label occupies its
box at the margin.

\noindent\rule{\textwidth}{0.3pt}
\begin{prop}
  \item[Physicalism, align=runin] Everything is physical, and this item is long
  enough to wrap onto a second line so that the continuation can be compared.
  \item[A considerably longer name, align=runin] Second item.
\end{prop}
\noindent\rule{\textwidth}{0.3pt}
\begin{prop}
  \item[Physicalism, align=flush] Everything is physical, and this item is long
  enough to wrap onto a second line so that the continuation can be compared.
  \item[A considerably longer name, align=flush] Second item.
\end{prop}
\noindent\rule{\textwidth}{0.3pt}
\begin{prop}
  \item[Physicalism] Everything is physical (the default, \texttt{left}).
\end{prop}
\noindent\rule{\textwidth}{0.3pt}

\section{Exactly one space, however the source is written}

The body has already begun by the time the label is emitted, so the space
following \verb|\item[...]| in the source would otherwise be typeset on top of
the one the label supplies.  All three items below must look identical.

\begin{prop}
  \item[One, align=runin] Text.
  \item[Two, align=runin]      Text.
  \item[Three, align=runin]
  Text.
\end{prop}

\section{An empty label leaves no trace}

\verb|counter=none| with no name gives an item with nothing to display; it must
produce neither a label nor a leading space.

\begin{prop}
  \item[counter=none, align=runin] Body with no label at all.
\end{prop}

\section{A body that opens with a sublist}

The label is real content, so the paragraph it starts is not empty and the
sublist simply follows it.  The item after the sublist must be unaffected.

\begin{prop}
  \item[Outer, align=runin]
  \begin{prop}
    \item Sub-item.
  \end{prop}
  \item[Next, align=runin] After the sublist.
\end{prop}

\section{Numbered items, and \texttt{inlineprop}}

A run-in label need not be a name:

\begin{prop}
  \item[align=runin] A numbered run-in item. \label{r:num}
\end{prop}

Reference: \ref{r:num}, which must be \texttt{(1)} --- \texttt{align} governs
placement only, never the label's content or its reference.

\verb|align| has no effect inside \verb|inlineprop|, as for every other value:
\begin{inlineprop}
  \item[Inline, align=runin] the body follows inline, as usual.
\end{inlineprop}

\end{document}
